\begin{styletable}

Синус
    & $\sin x$
    & \verb|\sin x|
\\

Косинус
    & $\cos x$
    & \verb|\cos x|
\\

Тангенс
    & $\tg x$
    & \verb|\tg x|
\\

Котангенс
    & $\ctg x$
    & \verb|\ctg x|
\\

Арксинус
    & $\arcsin x$
    & \verb|\arcsin x|
\\

Арккосинус
    & $\arccos x$
    & \verb|\arccos x|
\\

Арктангенс
    & $\arctg x$
    & \verb|\arctg x|
\\

Арккотангенс
    & $\arcctg x$
    & \verb|\arcctg x|
\\

Натуральный логарифм
    & $\ln x$
    & \verb|\ln x|
\\

Логарифм по основанию $a$
    & $\log_{a} x$
    & \verb|\log_{a} x|
\\

Экспонента
    & $\exp(x)$
    & \verb|\exp(x)|
\\

\hline
\end{styletable}

\stylenote{
    Скобки вокруг одиночного простого аргумента тригонометрических и логарифмических функций не обязательны ($\sin x$, $\cos \varphi$, $\tg \alpha, \ln x$, и~т.п.).
    Скобки ставятся, если аргумент~--- составное выражение или произведение, например: $\sin(x + y)$, $\cos(2 \pi n)$, $\tg(\alpha / 2)$.
}